\section{Past Sampling at Fixed Forecast Anchors}
All arms retain the same queries, future labels, CV and frozen
half-blend recipients. We fit the alternative velocity to offsets
$(-3,-2,-1,0)$, $(-6,-4,-2,0)$, $(-6,-5,-1,0)$, or every
offset from $-6$ to $0$. Slopes use native index units and paths
remain anchored at the latest position.

\begin{table}[!t]
\centering\small
\caption{Direct conservative gain vs the recent fit (points)
and original passes on 24 fitting / 40 opened recordings.}
\begin{tabular}{@{}lrrr@{}}
\toprule
History & Fit & Opened & Passes\\
\midrule
\PastSampleRows
\bottomrule
\end{tabular}
\end{table}

The 95\% paired source intervals for wide uniform versus recent
are \PastSampleSource{} on fitting and \PastSampleTarget{} on
opened recordings. This is sensitivity to past sampling, not a
claim of matched physical time or a transfer remedy.

We test an optimistic relaxation of the candidate family on fitting
labels only. At each selected row and known future step, distance
to the convex hull of the four paths lower-bounds any mixture's
error, even allowing the mixture to vary with the step. For wholly
missing rows, terminal CV-to-hull distance lower-bounds the maximum
disagreement debit. Taking positive harm above CV and using the
largest possible selected reference mass yields a risk lower bound.
This excludes \PastSampleBlocked/\PastSampleFit{} recordings.
It rules out reweighting these schedules at fixed recipients and
half-blend amplitudes, not different candidate paths or recipients.
